Im Kraftwerk wird Brennstoff mit dem Energieinhalt $Q$ verbrannt. Daraus entsteht die elektrische Energie $W = \eta\,Q$, mit der die Wärmepumpe die Wärme
\begin{align}
 Q_\text{w} &= \eta_\text{WP}\,W = \eta_\text{WP}\,\eta\,Q
\end{align}
ins Haus liefert. Würde man den Brennstoff direkt im Haus verbrennen, erhielte man (höchstens) die Wärme $Q$. Die Wärmepumpe lohnt sich also, wenn $\eta_\text{WP}\,\eta\,Q > Q$ ist, also
\begin{align}
 \eta_\text{WP} &> \ewpF = \frac{1}{\etKO} = \ewp \approx \result{\ewpP{0}{2}}
\end{align}
Verluste bei der Übertragung der elektrischen Energie und in der Ölheizung sind dabei vernachlässigt.
